%%%%

%%%   Самарский государственный технический университет

%%%   Кафедра прикладной математики и информатики

%%%

%%%   Преамбула для статьи в журнал "Вестник Самарского государственного

%%%   технического университета. Серия: «Физико-математические науки»"

%%%   Ориентирована под MikTEx 2.4 for Win32

%%%

%%%   Хотя сам журнал собирается под GNU/Linux на дистрибутиве TeXLive



\documentclass[11pt,a4paper,twoside]{article}



\usepackage[cp1251]{inputenc}

\usepackage[T2A]{fontenc}

\usepackage[english,russian]{babel}

\usepackage{samgtubib}

\usepackage[dvips]{graphicx}

\usepackage[multiple]{footmisc}



\usepackage{samgtu}

\renewcommand{\VestnikYear}{2009} % Год издания Вестника

\renewcommand{\VestnikNumber}{2\,(19)} % Номер Вестника

\renewcommand{\VestnikMonth}{Ноябрь} % Месяц выхода Вестника

\renewcommand{\VestnikData}{01.11.09} % Дата подписания Вестника в печать

\renewcommand{\VestnikUPL}{26,64} % Усл. печ. л.

\renewcommand{\VestnikUIL}{25,73} % Уч.-изд. л.

\renewcommand{\VestnikRegNumber}{479/09} % Регистрационный номер

\renewcommand{\VestnikOrder}{994} % Номер заказа






%
%\begin{document}

\renewcommand{\firstpage}{8}
\renewcommand{\lastpage}{16}


\udk{517.956.45}
\msc{35R35; 35K05, 35R05}

\title{Нелокальная задача Стефана для квазилинейного параболического уравнения}
{Нелокальная задача Стефана для квазилинейного параболического уравнения}

\author{Ж.\,О.~Тахиров\inst{1}, Р.\,Н.~Тураев\inst{2}}{Т\,а\,х\,и\,р\,о\,в~Ж.\,О.,\ Т\,у\,р\,а\,е\,в~Р.\,Н.}

\institute{Ташкентский государственный педагогический университет им.~Низами, \\
700100, Узбекистан,  Ташкент, ул. Ю.~Х.~Хажиба, 103. \and 
Институт математики и информационных технологий АН Республики Узбекистан,\\ 
100125, Узбекистан,  Ташкент, ул. Ф.~Ходжаева, 29.
}

\email{E-mails}{prof.takhirov@yahoo.com, rasul.turaev@mail.ru}

\otherinf{\fullauthor{Жозил Останович Тахиров} \regalia{д.ф.-м.н., проф.} \position{зав. кафедрой}
\dept{каф. математического анализа}
\fullauthor{Расул Нортожиевич Тураев} \regalia{к.ф.-м.н.} \position{старший научный сотрудник}
\dept{лаб. дифференциальных уравнений}}

\keywords{нелокальная задача, задача Стефана, квазилинейное параболическое уравнение, свободная граница,
априорные оценки, теорема существования и единственности, фиксированная граница, метод потенциалов, принцип максимума}

\workabstract{Рассматривается задача со свободной границей с нелокальным
граничным условием для квазилинейного гиперболического уравнения. Для искомого решения установлены априорные
оценки Шаудеровского типа. На основе полученных оценок доказаны
теоремы единственности и существования.}

\engtitle{The nonlocal Stefan problem for quasilinear parabolic equation}

\engauthor{J.\,O.~Takhirov\inst{1}, R.\,N.~Turaev\inst{2}}{T\,a\,k\,h\,i\,r\,o\,v~J.\,O.,\ T\,u\,r\,a\,e\,v~R.\,N.}

\enginstitute{Nizami Tashkent State Pedagogical University, \\
103, Yusuf Khos Khojib st., Tashkent, 700100, Uzbekistan.
\and
Institute of Mathematics and Information Technologies \\ of the Academy of Sciences of Uzbekistan, \\
29, Do'rmon yo'li srt., Tashkent, 100125,  Uzbekistan.}




\otherenginf{\fullauthor{Jozil O. Takhirov} \regalia{Dr.\,Sci. (Phys. \& Math.)} \position{Head of Dept.}
\dept{Dept. of Mathematical Analysis}
\fullauthor{Rasul N. Turaev} \regalia{Ph.\,D. (Phys. \& Math.)} \position{Senior Research Scientist}
\dept{Lab. of Differential Equations}}

\engabstract{In this paper, we deal with free  boundary problem with nonlocal boundary condition
for quasilinear parabolic equation. For the solutions of the problem apriory estimates
of Shauder's type are established. On the base of apriory estimations the existence
and uniqueness theorems are proved.}

\engkeywords{nonlocal problem, Stefan problem, quasilinear parabolic equation, free boundary,
priori estimates, existence and uniqueness theorem, fixed boundary, method of potentials, maximum principle}

\date{10/VII/2011}
\revisiondate{19/XII/2011}

\maketitle




\Section[n]{Введение}
Теория классической разрешимости задачи Стефана и других задач со
свободными границами для параболических уравнений построена в~работах А.~Фридмана \cite{takh:bib:2}, А.~Мейрманова \cite{takh:bib:3}, Л.~Рубинштейна \cite{takh:bib:1} и~др.

Особенность данной задачи состоит в наличии переменных размеров области, в~которой исследуется температурное поле, и наличия подвижной границы раздела фаз, изучение поведения которой и составляет основную цель решения. Физические свойства среды, находящейся в разных фазах, будут различными. Поэтому задача Стефана характеризуется существенной геометрической и физической
нелинейностью, что крайне затрудняет её решение.

В~работе \cite{takh:bib:4} Джим Дуглас  предлагает способ доказательства
единственности решения задачи Стефана для квазилинейного
параболического уравнения. В.~Кайнером \cite{takh:bib:5} были доказаны
теоремы существование и единственности однофазной задачи Стефана для нелинейных уравнений.

Отличительной особенностью однофазных задач является монотонность
свободной границы и, как следствие этого, ограниченность первой
производной решения по пространственной переменной на свободной
границе.

Цель настоящей статьи "--- вывод априорных оценок решений однофазной задачи Стефана с нелокальным  граничным условием на фиксированной границе, обеспечивающих её разрешимость для любого конечного промежутка времени, и~доказательство единственности решения.

\smallskip

\Section{Постановка задачи} Требуется найти пару функций
$(s(t), u(x, t))$ удовлетворяющих  условиям
\begin{gather}
a(u)u_{t}=u_{xx}, \quad \text{в} \quad  D=\{(x,t):0<x<s(t), \; 0<t\leq T \}; 
\label{takh:1}
\\
u(x, 0)=\varphi(x), \quad 0\leq x \leq s_0; 
\label{takh:3}
\\
u(0, t)=mu(x_0,t), \quad 0\leq t \leq T, \ 0< x_0 \leq s_0;
\label{takh:3} 
\\
u(s(t), t)=0,\quad 0\leq t \leq T; 
\label{takh:5}
\\
\dot{s}(t)=-cu_x(s(t),t),\quad 0\leq t \leq T. 
\label{takh:5}
\end{gather}


Будем предполагать выполнение следующих условий:
\begin{itemize}
\item[a)] функции $a(u)$ и $a'(u)$ определены для любого значения
аргумента и ограничены на любом замкнутом множестве аргумента,
причём $a(u)\hm \geq a_0>0$, $a'(u)>0$;

\item[b)] постоянные $s_0$, $m$, $x_0$ и $c$ удовлетворяют следующим неравенствам:
\mbox{$s_0{>}0$}, $0<m<1$, $c>0$, $0<x_0\leq s_0$;


\item[c)] $0\leq \varphi(x)\leq N(s_0-x)$, $N=\max\limits_{0\leq x\leq
s_0}|\varphi (x)(s_0-x)^{-1}|$.
\end{itemize}

Задача \eqref{takh:1}--\eqref{takh:5} исследованы в работах \cite{takh:bib:3, takh:bib:4} в случае, когда вместо \eqref{takh:3} задаётся граничное условие второго рода.

\smallskip

\Section{Некоторые априорные оценки}

\begin{lemma}\sl Пусть выполнены условия {\rm a), b), c).} Тогда для решения $u(x,t),$ $s(t)$ задачи \eqref{takh:1}--\eqref{takh:5} справедливы оценки
\begin{gather}
0\leq u(x,t)\leq M_1=\max\limits_{x}|\varphi(x)|,\quad  (x,t)\in
\bar{D},\label{takh:6}
\\
0<\dot{s}(t)\leq M_2=CN,\quad  0<t\leq T,\label{takh:7}
\\
0\leq u(x,t)\leq N(s(t)-x),\quad  0\leq x\leq s(t),\; 0\leq t\leq T.
\label{takh:8}
\end{gather}
\end{lemma}


\begin{proof} Функция $u(x,t)$ удовлетворяет однородному параболическому уравнению, для которого справедлив
принцип максимума, и ограничена на границе $x=s(t)$ и в начальный момент времени. 
В возможных точках положительного максимума или отрицательного минимума $u(t, 0)$  оценивается через
известные величины.

Неположительность производной $u_x(s(t),t)$ следует из положительности решения $u(x,t)$ в области $D$ и равенства его
нулю на свободной границе.

Для оценки производной $u_x(s(t),t)$ снизу рассмотрим функцию
\begin{equation}
v(x,t)=u(x,t)+N(x-s(t)),\quad  N={\rm const}>0.
\label{takh:9}
\end{equation}
Тогда из задачи \eqref{takh:1}--\eqref{takh:5} получим задачу для $v(x,t)$:
\begin{gather}
v_{xx}-av_t=aN\dot{s}(t)\geq 0,\quad  (x,t)\in D,\label{takh:10}
\\
v(0,t)=mv(x_0,t)-mNx_0-(1-m) Ns(t),\label{takh:11}
\\
v(x, 0)=\varphi(x)+N(x-s_0)\leq 0,\label{takh:12}
\\
v(t,s(t))=0.\label{takh:13}
\end{gather}

Из \eqref{takh:9} следует, что $v(x,t)$  не может достигать положительного максимума внутри области $D$. 
Условие \eqref{takh:11} не позволяет допустить существования на левой границе положительного максимума. 
Таким образом, функция $v(x,t)$ неположительна всюду в $D$ и~равна нулю на границе $x=s(t)$. Но тогда $v_x(s(t),t)\geq 0$. Следовательно,
\begin{equation}
v_x(s(t),t)=u_x(s(t),t)+N\geq 0
\label{takh:14}
\end{equation}
или
$$u_x(s(t),t)\geq -N.$$
Из неположительности $v(x,t)$ с учётом \eqref{takh:9} можно получить
неравенство \eqref{takh:8}.\end{proof}

\smallskip

Теперь, используя результаты \cite{takh:bib:6}, получим оценки для $|u_x|$ и~$|u|^{D}_{1+\gamma}$. 
Здесь и далее в отношении функциональных пространств и обозначения норм в них будем следовать обозначениям
работы \cite{takh:bib:6}.

Пусть
$Q=\{(x,t) : 0\leq x\leq s_0, 0\leq t\leq T\}$,
$Q^{\delta}=\{(x,t) : 0<\delta\leq x \hm\leq {s_0-\delta}, 0<\delta\leq t\leq T\}$,
$
Q^{\delta}_{0}=\{(x,t) : 0<\delta\leq x\leq s_0-\delta, 0\leq t\leq T\}$.

\smallskip


\begin{theorem}[1]\sl Пусть функция $u(x,t)$ ($M=\max\limits_{Q}|u|$
непрерывна в $Q$  вместе с производной $u_x$ и удовлетворяет
уравнению \eqref{takh:1} всюду в $Q$ за исключением, может быть, точек прямой
$x=0.$ Тогда
\begin{equation}
|u_x(x,t)|\leq P_0(M,a_0,\delta), (x,t)\in Q^{\delta}_{0}.
\label{takh:15}
\end{equation}
Если ещё известно$,$ что функция $u(x,t)$ обладает в  $Q$ суммируемыми с~квадратом обобщёнными производными $u_{xx}$ и
$u_{tx},$ то существует такое $\gamma =\gamma(M,a_0,\delta),$ что
\begin{equation}
|u|^{Q^{2\delta}_{0}}_{1+\gamma}\leq C(M,a_0,a_1,\delta),\quad  
0<\gamma<1,
\label{takh:16}
\end{equation}
где $a_1=\max a(u)$ при $|u|\leq M.$

Пусть в области $\{(x,t)\in Q, |u|\leq M \}$ функция $a(u)$ удовлетворяет условию Гёльдера с показателем $\beta$ по $u$ и с константой $N,$ а функция $u(x,t)$ непрерывна вместе с~производными $u_t,$ $u_x,$ $u_{xx}$, причём $u(x, 0)=0.$ 
Положим $\alpha=\beta\gamma$ и предположим{\rm ,} что $|u|^{Q^{\delta}_{0}}_{2+\gamma}<\infty .$
Тогда \begin{equation}
|u|^{Q^{4\delta}_{0}}_{2+\gamma}\leq C(M,a_0,a_1,N,\beta,\delta).
\label{takh:17}
\end{equation}
\end{theorem}

\smallskip


\begin{proof} Для  $(x,t)\in Q^{\delta}_{0}$ оценки
\eqref{takh:15}, \eqref{takh:16}, \eqref{takh:17} непосредственно следуют из результатов работы
\cite{takh:bib:6}. Теперь, выбирая $\delta$ таким, что $4\delta\leq x_0$, и
пользуясь нелокальным условием \eqref{takh:3} с учётом \eqref{takh:17}, получаем,
что ограниченная функция $u_t(x,t)$ удовлетворяет условию Гёльдера
с показателем $\alpha$. Оценим $|u_x|$ в области
$Q_1=\{(x,t):0\leq x\leq 4\delta, 0\leq t\leq T\}$. В $Q_1$ введём
новую функцию
\begin{equation}
v(x,t)=u(x,t)-g_1(x,t), 
\label{takh:18}
\end{equation}
где
$$
g_1(x,t)=\frac{4\delta-x}{4\delta}u(0,t)+\frac{x}{4\delta}u(4\delta ,t).
$$
Для функции $v(t,x)$ в области $Q_1$ имеем задачу
\begin{gather*}
a(v)v_t=v_{xx}+f(x,t);\quad 
v(x, 0)=0, \quad  v(0,t)=0, \quad  v(4\delta ,t)=0.
\end{gather*}
Так как $v(x,t)=0$ на $\Gamma$ $(\Gamma=\partial Q_1)$,  оценки
\eqref{takh:15}, \eqref{takh:16}, \eqref{takh:17} выполняются всюду в $Q_1$:
\begin{equation}
|u_x|\leq C, \quad  |u|^{Q_1}_{1+\gamma}\leq C,  \quad |u|^{Q_1}_{2+\gamma}\leq C,\quad (x,t)\in Q_1.
\label{takh:19}
\end{equation}
Чтобы получить оценки вблизи правой (свободной) границы, выполним
замену независимых переменных $\tau =t$, $y={x}/{s(t)}$. Тогда
области $D$ соответствует область $\Omega = \{(y,\tau ): 0<y<1,
0<\tau \leq T\}$, а ограниченная функция $v(y,\tau )=u(ys(\tau),\tau )$ является решением задачи
\begin{gather}
\bar{a}v_{\tau}=\frac{1}{s^2(\tau)}v_{yy}+f(y,\tau), \quad (y,\tau)\in \Omega;\label{takh:20}
\\
v(0,\tau)=mv(h(\tau),\tau), \quad %0\leq \tau \leq T,\label{takh:21} \\
v(y, 0)=\varphi y(s_0),\quad  %0\leq \tau \leq T,\label{takh:22} \\
v(1,\tau)=0, \quad 0\leq \tau \leq T,\label{takh:23}
\end{gather}
где
$$
f(y,\tau)=-\bar{a}\frac{y}{s^2(\tau)}v_y(y,\tau)v_y(1,\tau),
\quad h(\tau)={s_0}/{s(\tau)}>0.
$$
Коэффициенты и правая часть уравнения \eqref{takh:20} удовлетворяют условиям
Гёльдера и теоремы 1. Применяя метод четного продолжения через
правую границу \cite{takh:bib:6},  получим оценки для $|v_y|$, $|v|_{1+\gamma}$
вплоть до $y=1$. \end{proof}

\smallskip

Априорные оценки старших производных устанавливаются при помощи
результатов по линейным уравнениям.

\begin{theorem}[2]\sl Пусть коэффициенты уравнения
\begin{equation}
a(x,t)v_{xx}+b(x,t)v_x+c(x,t)v-v_t=f(x,t), (x,t)\in
\Omega
\label{takh:24}
\end{equation}
удовлетворяют условиям Гёльдера
\begin{gather*}
|a|^{\bar{\Omega}}_{\gamma}+|b|^{\bar{\Omega}}_{\gamma}+|c|^{\Omega}_{\gamma}+|f|^{\bar{\Omega}}_{\gamma}<\infty,
\quad 
a(x,t)\geq a_0>0.
\end{gather*}
Пусть $v(x,t)$ "--- решение уравнения \eqref{takh:24}{\rm ,} которое
 удовлетворяет условиям
\begin{equation}
v(x, 0)=\varphi(x),\quad v(0,t)=mv(h(t),t),\quad v(1, t)=0,
\label{takh:25}
\end{equation}
причём $|v|^{\Omega}_{2+\gamma}<+\infty,$ $M=\max\limits_{\Omega}|v(x,t)|.$
Тогда
\begin{equation}
|v|^{\bar{\Omega}}_{2+\gamma}\leq C(|f|^{\Omega}_{\gamma}+M).
\label{takh:26}
\end{equation}
\end{theorem}

\smallskip

\textit{Д\,о\,к\,а\,з\,а\,т\,е\,л\,ь\,с\,т\,в\,о\/} проводится по следующей схеме. При получении
априорных оценок на левой половине области применяется способ,
который был использован при доказательстве теоремы 1. На правой
половине используем результаты работы \cite{takh:bib:2}. Разрешимость задачи
\eqref{takh:24}, \eqref{takh:25} доказывается методом потенциалов.

\smallskip

\Section{Единственность решения}
Получим интегральное представление для~$s(t)$, эквивалентное \eqref{takh:5}, для этого умножим уравнение \eqref{takh:1} на ($-x$) и проинтегрируем  полученное
выражение по области $D$:
\begin{equation}
\iint_{D}[(u-\xi u_{x})_x+(\xi b(u))_t]d\xi d\eta =\int_{\Gamma}(\xi-\xi u_x)d\eta-\xi b(u)d\xi=0,
\label{takh:27}
\end{equation}
$$
b(u)=\int^{u}_{0}a(\xi)d\xi.
$$
Имеем
\begin{equation}
s^2(t)=s^{2}_{0}+2\int^{s_0}_{0}\xi b(\eta)d\xi-2\int^{s(t)}_{0}\xi
b(u(\xi,t))d\xi+2\int^{t}_{0}u(0,\eta)d\eta.
\label{takh:28}
\end{equation}

\begin{theorem}[3]\sl При выполнении условий теоремы 1 решение
задачи \eqref{takh:1}--\eqref{takh:5} единственно{\rm .}
\end{theorem}

\begin{proof} Докажем  теорему для некоторого малого $t$, а затем покажем, что она справедлива для любого
$0<t<\infty$.

Пусть функции $s_1(t)$, $u_1(x,t)$ и $s_2(t)$, $u_2(x,t)$ являются
решениями задачи \eqref{takh:1}--\eqref{takh:5}. Пусть
$$
y(t)=\min(s_1(t),s_2(t)), \quad z(t)=\max(s_1(t),s_2(t)).
$$
Для каждой пары справедливо представление \eqref{takh:28}. Имеем
$$
s_1^2(t)-s_2^2(t)=2\int^{s_1(t)}_{0}\xi b(u_1)d\xi+2\int^{s_2(t)}_{0}\xi b(u_2)d\xi
+2\int^{t}_{0}(u_1(0,\eta)-u_2(0,\eta))d\eta.
$$
Отсюда
\begin{multline}
|s_1(t)-s_2(t)|\leq\frac{1}{s_0}\int^{y(t)}_{0}\xi|b(u_1)-b(u_2)|d\xi
+ \frac{1}{s_0}\int^{z(t)}_{y(t)}\xi b(u_i)d\xi +
\\
+\frac{1}{s_0}\int^{t}_{0}|u_1(0,\eta)-u_2(0,\eta)|d\eta,
\label{takh:29}
\end{multline}
где $u_i$ "--- решение между $y(t)$ и $z(t)$.

По лемме 
\begin{equation}
|u_i(x,t)|\leq N(y(t)-x),
\quad 
|u_1(y(t),t)-u_2(y(t),t)|\leq N|s_1(t)-s_2(t)|.
\label{takh:30}
\end{equation}
Рассмотрим функцию  $v(x,t)=u_1(x,t)-u_2(x,t)$. Тогда для $v(x,t)$
получим уравнение с ограниченными коэффициентами и следующую задачу
\begin{gather}
v_{xx}=b_1(x,t)v_t+b_2(x,t)v;
\\
%\begin{equation}
v(0,t)=mv(x_0,t), \quad  v(x, 0)=0,\\
%\label{takh:31}
%\end{equation}
%$$
|v(s(t),t)|\leq N|s_1(t)-s_2(t)|, 
\end{gather}
где $b_1(x,t)=a(u_1(x,t))$, $b_2(x,t)=[b(u_1)-b(u_2)]$.
Отсюда по
принципу максимума $$|u_1(x,t)-u_2(x,t)|\leq N\max\limits_{0\leq
\eta \leq t}|s_1(\eta)-s_2(\eta)|.$$

В силу ограниченности функций $u(x,t)$, $a(u)$, $a'(u)$ оценим члены из
\eqref{takh:29}:
\begin{multline}
I_1=\frac{1}{s_0}\int_0^{y(t)}\xi|b(u_1)-b(u_2)|d\xi\leq\frac{b^{*}(\overline{\xi})}{s_0}
y^2(t)|u_1-u_2|\leq
\\
\leq M_{11}\max\limits_{\eta}|s_1(\eta)-s_2(\eta)|y^2(t),
\label{takh:32}
\end{multline}
\begin{multline}
I_2=\frac{1}{s_0}\int_{y(t)}^{z(t)}\xi b(u_i)d\xi=
\frac{1}{s_0}\int_{y(t)}^{z(t)}\xi
d\xi\int_{0}^{u_i}a(\chi) d\chi\leq
\frac{a_1}{s_0}\int_{y(t)}^{z(t)}\xi u_i(\xi,t)d\xi\leq
\\
\leq\frac{a_1}{s_0}\int_{y(t)}^{z(t)}\xi N(z(t)-\xi)d\xi=
\frac{a_1N}{s_0}\int_{y(t)}^{z(t)}\xi (z(t)-\xi)d\xi=
\\
=\frac{-M_{12}}{2}\int_{y(t)}^{z(t)}\xi d(z(t)-\xi)^2=
\frac{M_{12}}{2}y(t)(z(t)-y(t))^2+
\\
+M_{12}\int_{y(t)}^{z(t)}(z(t)-\xi)^2d\xi=
M_{12}y(t)(z(t)-y(t))^2+M_{12}(z(t)-y(t))^3\leq
\\
\leq M_{12}y(t)\max\limits_{0\leq\eta\leq
t}|s_1(\eta)-s_2(\eta)|^2+ M_{12}\max\limits_{0\leq\eta\leq
t}|s_1(\eta)-s_2(\eta)|^3,
\label{takh:33}
\end{multline}
\begin{equation}
I_3=\frac{1}{s_0}\int_0^t|u_1(0,\eta)-u_2(0,\eta)|d\eta\leq
M_{13}T\max\limits_{0\leq\eta\leq
t}|s_1(\eta)-s_2(\eta)|.
\label{takh:34}
\end{equation}
Пусть $A(T)=\max\limits_{0\leq t\leq T}|s_1(t)-s_2(t)|$. Если $A(T)>0$, то с учётом полученных
неравенств из \eqref{takh:29} имеем
\begin{equation}
A(T)\leq
M_{11}y^2(T)A(T)+M_{12}y(t)A^2(T)+M_{12}A^3(T)+M_{13}T\cdot A(T),
\label{takh:35}
\end{equation}
Если \eqref{takh:35} разделить на $A(T)$, то получим
$$
1\leq M_{11}y^2(T)+M_{12}y(T)A(T)+M_{12}A^2(T)+M_3T.
$$

Получим сформулированное выше утверждение для малых $t$. 
Пусть для определённости $T<1$. Тогда вместо $T^2$ берём $T$:
\begin{gather*}
1\leq M_{11}T_0+M_{12}T_0+M_{13}T_0,
\quad 
1\leq T_0(M_{11}+M_{12}+M_{13})=T_0M_{14}.
\end{gather*}
Если $M_{14}T_0<1,$ то получим противоречие. Если $1\leq T,$ то
$1\leq T^2M_{15}$ и $T^2M_{15}<1$. Тогда для $0<t\leq T_0$
имеет место теорема единственности. 

Теперь докажем, что единственность имеет место для любого $0<t<+\infty$.
Пусть $T_1=\sup\limits\{t:s_1(\eta)=s_2(\eta), 0\leq \eta\leq
t\}$. Если $T_1=\infty$, то доказательство очевидно. Предположим, что величина
$T_1$ ограничена, и попробуем получить противоречие.
Пусть
$$
A(\Delta(t))=\max\limits_{T_1\leq t\leq T_1+\Delta
t}|s_1(t)-s_2(t)|.
$$
В силу \eqref{takh:7} имеем неравенство $A(\Delta t)\leq M_2\Delta t$. По предположению, выполняются неравенства
$A(\Delta t)>0$, $\Delta t>0$. Оценим члены неравенства \eqref{takh:29}:
$$
|I_2|=\left|\frac{1}{s_0}\int_{y(t)}^{z(t)}\xi
b(u_i)d\xi\right|\leq M_{17}A^2(\Delta t), \quad  T_1\leq t\leq
T_1+\Delta t.
$$



\begin{multline*}
I_3=\frac{1}{s_0}\int_{T_1}^{t}\left|u_1(0,\eta)-u_2(0,\eta)\right|d\eta= %\\=
\frac{m}{s_0}\int_{T_1}^{t}\left|u_1(x_0,\eta)-u_2(x_0,\eta)d\eta\right|\leq 
\\ \leq
\frac{M}{s_0}\Delta t\max\limits_{t}|s_1(t)-s_2(t)| %\leq \\
\leq\frac{M}{s_0}\Delta tA(\Delta t) \leq M_{18}A(\Delta t)\Delta t, \quad
 T_1\leq t\leq T_1+\Delta t.
\end{multline*}

\noindent
Для члена $I_1$, как и в первом случае, имеем
$$
A(\Delta t)\leq \!\max \limits_{T_1\leq t\leq T_1+\Delta t}M_{14} \!\! \int
_{0}^{y(t)} \!\!\! x|bu_1(x,t)-b_2(u(x,t)|dx+M_{17}A^2(\Delta
t)+M_{18}A(\Delta t)\Delta t.
$$
Отсюда
\begin{equation}
1\leq M_{15}\max\limits_{T_1\leq t\leq T_1+\Delta t}\int
_{0}^{y(t)}\frac{x|bu_1(x,t)-bu_2(x,t)|}{A(\Delta
t)}dx+M_{17}A(\Delta t) +M_{18}\Delta t. 
\label{takh:36}
\end{equation}
Перепишем \eqref{takh:36} в виде
\begin{multline}
1\leq M_{20}\max\limits_{t}\int
_{0}^{y(T_1)}\frac{x|u_1(x,t)-u_2(u(x,t)|}{A(\Delta t)}dx+
\\
+M_{20}\max\limits_{T_1\leq t\leq T_1+\Delta t}\int
_{y(t_1)}^{y(t)}\frac{x|u_1(x,t)-u_2(u(x,t)|}{A(\Delta
t)}dx+ M_{17}A(\Delta t)+M_{18}\Delta t
\label{takh:37}
\end{multline}
и рассмотрим функцию $v(x,t)=u_1(x,t)-u_2(x,t)$. Тогда имеем задачу
\begin{gather*}
v_{xx}=b_1(x,t)v_t+b_2(x,t)v,\quad  0\leq x\leq y(t), \quad  T_1\leq t\leq T_1+\Delta t;
\\
v(x,T_1)\equiv0, \quad  v(y(t),t)\leq NA(\Delta t), \quad  v(0, t)=mv(x_0, t).
\end{gather*}
Отсюда по принципу максимума
\begin{equation}
|v(x,t)|\leq NA(\Delta t). 
\label{takh:38}
\end{equation}
Оценим интегральные члены в формуле \eqref{takh:37}:
\begin{multline*}
\int_{y(T_1)}^{y(t)}x\frac{|v(x,t)|}{A(\Delta t)}dx\leq
\int_{y(T_1)}^{y(t)}xNdx = %\leq
%\frac{N}{2}x^2\left|^{y(t)}_{y(T_1)}\right.=
\frac{N}{2}\left(y^2(t)-y^2(T_1)\right)=
\\
=\frac{N}{2}\left(y(t)+y(T_1)\right)\left(y(t)-y(T_1)\right)\leq
M_{21}(t-T_1)\leq M_{21}\Delta t.
\end{multline*}

Теперь рассмотрим вспомогательную задачу:
\begin{gather*}
w_{xx}=b_1(x,t)w_t+b_2(x,t)w,\quad  0\leq x\leq y(T_1), \quad  T_1< t;
\\
w(x,T_1)=0, \quad   0<x<y(T_1); \\ w(0,t)=mw(x_0,t), \quad w(y(T_1),t)=1, \quad  t>T_1.
\end{gather*}
Введём функцию
$$
V(x,t)=\frac{v(x,t)}{NA(\Delta t)}-w(x,t), \quad  0<x<y(T_1), \quad 
T_1\leq t\leq T_1+\Delta t,
$$
для которой поставим такую задачу:
\begin{gather*}
V_{xx}=b_1V_t+b_2V, \quad  0<x<y(T_1), \quad  T_1<t;
\\
V(x,T_1)=0, \quad  V(0, t)=mV(x_0, t), \quad  T_1<t;
\\
V(y(T_1),t)=\frac{v(y(T_1),t)}{NA(\Delta t)}-w(y(T_1),t)\leq 0.
\end{gather*}
По принципу максимума $V(x,t)\leq0$, следовательно,
$$
\frac{|v(x,t)|}{NA(\Delta t)}\leq w(x,t)\leq1, \ \ 0\leq x\leq
y(T_1), \ \ T_1\leq t\leq T_1+\Delta t.
$$
В промежутке $0\leq x\leq y(T_1)$ имеет место
$$
\lim\limits_{t\to  T_1}w(x,t)=0.
$$
Тогда 
$$
\lim\limits_{t\to  T_1}\int_{0}^{y(T_1)}w(x,t)dx=0,
$$
следовательно,
$$
\max\limits_{T_1\leq t\leq T_1+\Delta t}\int_{0}^{y(T_1)}
x\frac{|v(x,t)|}{NA(\Delta t)}dx\to 0, \ \ \text{при }
\Delta t\to 0.
$$
Если в формуле \eqref{takh:38} перейти к пределу при $\Delta t\to 0$,
то получим противоречие.
Таким образом $s_1(t)\equiv s_2(t)$ для любого $t>0$. \end{proof}

\smallskip

\Section{Существование решения}
При определении максимального интервала существования решения
задач Стефана учитываются три фактора: 
\begin{enumerate}
\item невырожденность области;
\item  наличие априорных оценок норм в соответствующем пространстве;
\item  ограниченность снизу и сверху модуля градиента решения на
свободной границе.
\end{enumerate}

Если наложить некоторые ограничения  (обеспечивающие выполнение вышеуказанных факторов на произвольном интервале времени) на данные задачи, то классическое решение задачи Стефана существует при всех положительных значениях времени.

\begin{theorem}[4]\sl При условиях теорем 1 и 2{\rm ,} а также
соответствующих условиях согласования в точках $(0, 0),$ $(s_0, 0)$
существует единственное решение $u(x,t)\in C^{2+\gamma}(\overline{D}),$ $s(t)\in C^{1+\gamma}[0,T]$ задачи \rm
\eqref{takh:1}--\eqref{takh:5}.
\end{theorem}

\begin{proof} При доказательстве леммы  и теорем 1, 2 уже установлены необходимые априорные оценки. 
Свободная граница $x=s(t)$ монотонно возрастает с~ростом времени. Если гёльдеровость $\dot{s}(t)$ доказана и априорные оценки норм в пространстве $C^{2+\gamma}$ для $u(x,t)$ получены, то можно доказать глобальную разрешимость задачи. \end{proof}



%\begin{enumerate}
%\item {\it Рубинштейн Л.И.} Проблема Стефана [Текст]/Л.И.Рубинштейн.--Рига: Звайгзне, 1967.--468 с.
%\item {\it Фридман А.} Уравнения с частными производными
%параболического типа [Текст]/ А.Фридман.--М.: 1968.--428 с.
%\item {\it Мейрманов А.М.} Задача Стефана [Текст]/А.М.Мейрманов.--Новосибирск: Наука, 1986.--239 с.
%\item {\it J. Douglas,Jr.}  A uniqueness theorem for the solution of the Stefan problem [Текст]/J. Douglas,Jr.//
%Proc. Amer. Math. Soc.--1952.--Vol.8.--P. 402-408.
%\item {\it Kyner W.T.}  An existence and uniqueness theorem for a nonlinear Stefan problem [Текст]/ Kyner W.T.//
%J.Math.and Mech.--1959.--Vol.8 \No4.--P.483-498.
%\item  {\it Кружков С.Н.} Нелинейные параболические уравнения с двумя независимыми пе-ременными[Текст]/
%С.Н.Кружков// Труды Моск. Матем. Общ-ва.--1967.--Т.16.--
%С.329-346.
%\end{enumerate}

\begin{thebibliography}{9}

\Bibitem{takh:bib:2}
\book Partial Differential Equations of Parabolic Type
\by Friedman~A.
%MR181836 (31 #6062) 35.00 (35.62)
%Zbl 0144.34903
\publ Prentice-Hal
\publaddr Englewood Cliffs, N.J.
\yr  1964 
\totalpages xiv+347
\rtransl русск. пер. 
\by Фридман~А.
\book Уравнения с частными производными параболического типа
\publaddr М.
\publ Мир
\yr 1968
\totalpages 428

\RBibitem{takh:bib:3}
\by Мейрманов~А.\,М.
\book Задача Стефана
\publaddr Новосибирск
\publ Наука, Сибирск. отдел.
\yr 1986
\totalpages 240
\misctransl 
\by Meyrmanov~A.\,M. 
\book The Stefan problem
\publ Nauka, Sibirsk. Otdel.
\publaddr Novosibirsk
\yr 1986
\totalpages 240


\RBibitem{takh:bib:1}
\by Рубинштейн~Л.\,И.
\book Проблема Стефана
\publaddr Рига
\publ Звайгзне
\yr 1967
\totalpages 468
\misctransl
\by Rubinshteyn~L.\,I.
\book The Stefan problem
\publ Zvaygzne
\publaddr Riga 
\yr 1967 
\totalpages 468

\Bibitem{takh:bib:4}
\by Douglas~J., Jr.
\paper A uniqueness theorem for the solution of the Stefan problem
\jour Proc. Amer. Math. Soc.
\yr 1952
\vol 8
\pages 402--408




\Bibitem{takh:bib:5}
\by Kyner~W.\,T.
\paper An existence and uniqueness theorem for a nonlinear Stefan problem
\jour J. Math. and Mech.
\yr 1959
\vol 8
\issue 4
\pages 483--498



\RBibitem{takh:bib:6}
\by Кружков~С.\,Н.
\paper Нелинейные параболические уравнения с двумя независимыми переменными
\jour Тр. Москов. матем. об-ва
\yr 1967
\vol 16
\pages 329--346
\misctransl
\by Kruzhkov~S.\,N.
\paper Nonlinear parabolic equations in two independent variables
\jour Trudy Moskov. Matem. Obshch.
\yr 1967
\vol 16
\pages 329--346


\end{thebibliography}
 
\noindent
\hskip6.5cm \thedate \par\noindent
\hskip6.5cm\therevdate

\vspace{-5mm}

\makeengtitle



 \newpage
%
% \tableofengcontents
%
%\end{document}
